\chapter{Tables}
\label{chap:tables}

\section{A simple table}

\Cref{tab:simple} uses the \texttt{booktabs} rules (\verb|\toprule|,
\verb|\midrule|, \verb|\bottomrule|), which look cleaner than vertical lines.

\begin{table}[H]
  \centering
  \begin{tabular}{lrr}
    \toprule
    \textbf{Model} & \textbf{Accuracy (\%)} & \textbf{Time (s)} \\
    \midrule
    Baseline     & 71.2 & 0.8 \\
    Improved     & 84.5 & 1.4 \\
    Final        & \textbf{89.3} & 1.1 \\
    \bottomrule
  \end{tabular}
  \caption{A simple table with aligned numbers.}
  \label{tab:simple}
\end{table}

\section{A table as wide as the text}

With \texttt{tabularx}, the \texttt{X} columns share the remaining width and
wrap long text, as in \cref{tab:wide}.

\begin{table}[H]
  \centering
  \begin{tabularx}{\linewidth}{l X l}
    \toprule
    \textbf{Tool} & \textbf{Role} & \textbf{Needed by} \\
    \midrule
    \TeX{} distribution & Compilers and packages: MiKTeX on Windows, \TeX{} Live on Linux & everything \\
    Perl & Runs \texttt{latexmk}, which decides how many passes are needed & automatic builds \\
    Pygments & Highlights the code blocks of \cref{chap:code-math} & \texttt{minted} \\
    Biber & Sorts and formats the bibliography & \texttt{biblatex} \\
    \bottomrule
  \end{tabularx}
  \caption{A table that fills the text width.}
  \label{tab:wide}
\end{table}

\section{Merged cells and colors}

\Cref{tab:merged} merges cells across columns with \verb|\multicolumn| and
across rows with \verb|\multirow|, and colors a header and alternate rows.

\begin{table}[H]
  \centering
  \rowcolors{3}{gray!10}{white}
  \begin{tabular}{l l c c c}
    \toprule
    \rowcolor{blue!15}
    & & \multicolumn{3}{c}{\textbf{Score per sprint}} \\
    \rowcolor{blue!15}
    \textbf{Module} & \textbf{Task} & \textbf{S1} & \textbf{S2} & \textbf{S3} \\
    \midrule
    % With colored rows, put \multirow in the LAST row of the group, with a
    % negative count, or the next row's color hides its text
                               & Indexing   & 3 & 5 & 8 \\
    \multirow{-2}{*}{Search}   & Ranking    & 2 & 4 & 7 \\
                               & Components & 4 & 6 & 6 \\
    \multirow{-2}{*}{Frontend} & Tests      & 1 & 3 & 5 \\
    Backend                   & API        & 5 & 7 & 9 \\
    \bottomrule
  \end{tabular}
  \caption{Merged cells and colored rows.}
  \label{tab:merged}
\end{table}

\section{A table across several pages}

A \texttt{longtable} breaks across pages and repeats its header on each one.
It is not a floating environment, so it sits exactly where it is written
(\cref{tab:long}).

\begin{longtable}{c l l r}
  \caption{A long table that can span several pages.}
  \label{tab:long} \\
  \toprule
  \textbf{\#} & \textbf{Item} & \textbf{Category} & \textbf{Value} \\
  \midrule
  \endfirsthead
  \multicolumn{4}{l}{\small\itshape Continued from the previous page} \\
  \toprule
  \textbf{\#} & \textbf{Item} & \textbf{Category} & \textbf{Value} \\
  \midrule
  \endhead
  \midrule
  \multicolumn{4}{r}{\small\itshape Continued on the next page} \\
  \endfoot
  \bottomrule
  \endlastfoot
  1  & Requirements gathering & Planning     & 12 \\
  2  & State of the art       & Research     & 20 \\
  3  & Architecture           & Design       & 15 \\
  4  & Data collection        & Data         & 18 \\
  5  & Data cleaning          & Data         & 10 \\
  6  & First prototype        & Development  & 25 \\
  7  & Unit tests             & Quality      & 8  \\
  8  & Integration            & Development  & 14 \\
  9  & Performance tuning     & Development  & 11 \\
  10 & User feedback          & Validation   & 6  \\
  11 & Second iteration       & Development  & 22 \\
  12 & Documentation          & Writing      & 9  \\
  13 & Deployment             & Operations   & 7  \\
  14 & Monitoring             & Operations   & 5  \\
  15 & Report writing         & Writing      & 30 \\
  16 & Slides                 & Writing      & 6  \\
  17 & Rehearsal              & Presentation & 3  \\
  18 & Defense                & Presentation & 1  \\
\end{longtable}
